\section{Introduction}
\label{sec:intro}

Language-model agents now change the world rather than only describe it. They
write files, update databases, send messages, and run code, and most of this
work flows through tool servers written by third parties. The Model Context
Protocol (MCP) has become a common way to connect an agent to such servers. A
server advertises each tool with a name, a description, and an input schema; a
client calls the tool with arguments and receives content in
return~\cite{mcp2025tools}. The user, or a policy acting for the user, decides
whether a call is acceptable by reading that declaration and those arguments.

This arrangement has a gap that is easy to overlook. Approval is attached to
the declaration and the request, but the effect is produced by the server's
implementation, and the evidence of what happened is a response that the same
implementation writes. Consider a user who approves
\texttt{write\_file(path="report.txt", content="Q3 totals")}. A server whose
package was replaced after it was approved can write attacker bytes to
\texttt{report.txt}, or write the approved bytes somewhere else, and then
return \emph{``Successfully wrote to report.txt''}, the same message an honest
server returns. Schema validation of the response, a language-model judge, and
a signed receipt from the server all see identical evidence in the honest and
the malicious case. Authenticating the server tells the client who spoke. It
does not tell the client what that server did.

This paper asks whether a client can learn and enforce what an untrusted server
actually did, without trusting its testimony and without modifying the server.
We answer in two steps.

\paragraph{Responses are not evidence of effects.} We first make the gap
precise. If an honest implementation and a malicious one induce the same
distribution of requests and responses, then every monitor whose input is that
transcript has a detection rate against the malicious server equal to its
false-positive rate on the honest one (\S\ref{sec:measure}). The argument is a
standard indistinguishability step, but it has a practical consequence: more
careful response checking cannot help once the server keeps its story
consistent. Clients can do more than read one response, so we also built an
active auditor that issues follow-up reads and checks metamorphic
relations~\cite{chen2018metamorphic}, and we ran it across a funnel of 8,692
public registry servers. No configuration we tested reached the pre-registered
target of 20\% detection at no more than 10\% false positives. Strict settings
detected at most 0.5\% of attacks; permissive settings detected at most 7.0\%
while falsely flagging more than 70\% of honest servers.

\paragraph{Admit effects, do not audit testimony.} The negative result changes
the question. Instead of asking the server to describe its effect, the client
can decide which effects are allowed to reach state it trusts. MCPGate, the
system we build, runs the unmodified server against a private copy of the
relevant state. When the call returns, MCPGate stops every process that could
still write, reads the complete private result exactly once, compares it with a
contract derived from the approved call, and then either promotes those same
bytes into the trusted store or discards them. The server can do anything it
likes inside its private copy. What it cannot do is make an unapproved effect
appear in trusted state. Staging and commit are old ideas, and we build on
isolation systems such as Alcatraz~\cite{liang2009alcatraz} and
TxOS~\cite{porter2009txos}. Our contribution is to bind that mechanism to the
authorization object of an agent protocol and to show, on real servers, which
checks are needed for the binding to hold.

\paragraph{Evaluation.} We evaluate MCPGate against the strongest adversary our
threat model allows: a server that receives the approved arguments, returns
the honest response, performs a different effect, and also falsifies its own
read and listing results so that a response auditor sees a consistent honest
world. We realize this adversary on seven unmodified third-party servers by
interposing on the runtime libraries they use for file and database writes, so
the request and response stay honest and only the effect changes. The servers
were chosen and frozen before any security outcome was observed. Across
filesystem and SQLite servers in Node and Python, MCPGate prevented all 41
applicable attacks and blocked none of the honest workflows. A task-specific
static permission policy, which limits where a server may write but not what
it writes, prevented 20. A plain sandbox, the response auditor, and no defense
each prevented only 9, all of them attacks that happened to have no effect on
these servers. Per-call latency, CPU, and memory under MCPGate were within noise
of running without any defense.

\paragraph{Contributions.} This paper makes four contributions.
\begin{enumerate}
\item \textbf{An observation boundary with measurement.} We state when a
      response-only monitor cannot distinguish a diverted effect, and we measure
      active response auditing over a registry-scale funnel, reporting the
      selection effects and our own instrument failures
      (\S\ref{sec:measure}).
\item \textbf{A contract model for effects.} We separate destination,
      structure, and exact content, and we treat content that cannot be
      verified as an explicit UNKNOWN rather than silently granting it. Contracts
      for multi-file and database workflows are derived from honest runs of the
      approved version instead of being written by hand (\S\ref{sec:design}).
\item \textbf{An admission mechanism.} MCPGate combines request-shape binding,
      an atomic per-call allowance, private execution, writer shutdown, a single
      complete read, and promotion of the bytes that were checked
      (\S\ref{sec:design}, \S\ref{sec:impl}).
\item \textbf{A matched evaluation on frozen real servers.} We compare five
      conditions under one execution setup, in two effect domains, with
      server-clustered intervals, ablations, 100-trial race tests, and overhead
      (\S\ref{sec:method}, \S\ref{sec:results}).
\end{enumerate}

\paragraph{Scope.} MCPGate protects what enters a trusted store. It does not
prove that nothing happened elsewhere: a server with network access can still
send data out before MCPGate refuses its file result, and we observed a diverted
write land in a container-local path that MCPGate never promoted but also did
not prevent. MCPGate also assumes that the contract itself reflects what the
user wanted; it does not defend against a prompt that tricks the agent into
requesting the wrong thing. We return to these limits in
\S\ref{sec:limits}.
